\section{Differential Fuzzing}
\label{sec:greybox-differential}

The goal of differential fuzzing in \greyboxfuzz is to determine whether the target \zksnark program returns the correct verification result for a specific witness. Traditional differential fuzzing usually compares multiple implementations of the same specification. This strategy is difficult to apply directly to \zksnark programs because each application may encode a different statement, and building an independent \zksnark implementation for every statement is expensive and error-prone. Instead of requiring a second \zksnark implementation as the reference, \greyboxfuzz uses the intended statement itself as the test vector.

For each generated seed, \greyboxfuzz first evaluates whether the seed satisfies the statement that the target \zksnark program is intended to prove, and uses this evaluation to derive the expected verification result, namely accept or reject. The target prover and verifier are then executed on the same seed, and \greyboxfuzz compares the verifier's final decision against this expected result. Thus, the differential check is performed on the verification outcome only, rather than on the generated proof object itself. This distinction is important because two correct prover executions on the same input may still produce different proofs due to prover randomness, even though both proofs should lead to the same verification result. If the expected result and the verifier's decision agree, the seed does not expose an inconsistency. If they disagree, \greyboxfuzz reports a potential \logicerror. For example, if the statement is knowledge of a secret exponent $x$ such that $h = g^x$, where $g$ and $h$ are public values and $x$ is private, \greyboxfuzz does not require a second \zksnark prover and verifier as the reference. Instead, it only requires an independent program that checks whether the relation $h = g^x$ holds for the generated input and returns the corresponding accept-or-reject ground truth, without comparing the actual proofs themselves.

This design changes the role of test vectors in \greyboxfuzz. Rather than using fixed \zksnark test vectors or manually constructing a separate reference implementation, \greyboxfuzz generates new seeds during fuzzing and evaluates each seed against the statement. This is important because \logicerror often appears only for specific boundary cases, malformed but parseable values, or witness that violate the intended statement while still satisfying the program's expected format. By computing the expected result for each generated seed, \greyboxfuzz can detect cases where the target program accepts an invalid witness, rejects a valid witness, or produces a verification result inconsistent with the intended statement.

To reduce the burden of writing the \zksnark statement, \greyboxfuzz uses Charm \cite{akinyele2013charm}. Charm is a framework for prototyping cryptographic systems and provides high-level abstractions for common cryptographic objects and operations, including group elements, finite-field values, exponentiation, hashing, and pairing-based operations. Instead of asking the user to manually implement low-level cryptographic computations, \greyboxfuzz allows user to express the intended statement in Charm. The Charm program serves as the test vector and ground truth that computes the expected verification result for each generated seed.

For example, the statement $h = g^x$ can be expressed at a high level by declaring $g$ and $h$ as public group elements and $x$ as a private exponent. Conceptually, the statement has the following form with Charm:

\begin{ZZlisting}[H]
\centering
\begin{adjustbox}{max width=\textwidth}
\begin{CenteredBox}
\begin{minipage}{0.96\textwidth}
\begin{lstlisting}[style=pythonstyle]
from charm.toolbox.eccurve import prime192v1
from charm.toolbox.ecgroup import ECGroup

group = ECGroup(prime192v1)
g = group.random()
x = group.random()
h = g ** x

def oracle(g, h, x):
    return h == g ** x

\end{lstlisting}
\end{minipage}
\end{CenteredBox}
\end{adjustbox}
\caption{An example of Charm implementation}
\label{list:charm}
\end{ZZlisting}


The actual Charm implementation handles the underlying cryptographic representation of the variables, so the user does not need to manually implement elliptic-curve arithmetic or finite-field operations. This makes the statement easier to write and less likely to introduce unrelated implementation mistakes. The user only needs to describe the mathematical relationship that the proof is intended to enforce.

After the statement is defined, \greyboxfuzz uses it during fuzzing as follows. First, the input formatter generates a seed according to the expected structure of the statement. Second, the seed is evaluated by Charm to obtain the expected accept or reject result. Third, the same seed is passed to the target \pprog and \vprog to obtain the actual verification result. Finally, \greyboxfuzz compares the expected result with the actual result. A mismatch indicates that the target \zksnark program may not correctly implement the intended statement.

This statement-level test vector provides two advantages. First, it avoids the need to construct a separate \zksnark implementation for every target program. Second, it allows \greyboxfuzz to test the semantic correctness of the target program rather than only its ability to execute without crashing. As a result, \greyboxfuzz can detect \logicerror that would be missed by crash-oriented fuzzing, including missing constraints, incorrect witness conversion, invalid boundary handling, and incorrect proof or verification behavior.


